\documentclass[11pt]{article}
\usepackage{amsmath,amsthm,amssymb}
\usepackage[margin=1in]{geometry}
\usepackage{hyperref}
\usepackage{enumitem}
\usepackage{booktabs}

\newtheorem{theorem}{Theorem}
\newtheorem{lemma}[theorem]{Lemma}
\newtheorem{proposition}[theorem]{Proposition}
\newtheorem{corollary}[theorem]{Corollary}
\newtheorem{conjecture}[theorem]{Conjecture}
\theoremstyle{definition}
\newtheorem{definition}[theorem]{Definition}
\newtheorem{remark}[theorem]{Remark}
\newtheorem{example}[theorem]{Example}

\title{A recursive dimension formula for $B_{\ell+1}^{(\lambda)} V_\lambda$\\
  at multi-corner partitions}
\author{Clio}
\date{13 May 2026}

\begin{document}
\maketitle

\begin{abstract}
Let $H_q(S_n)$ be the Iwahori--Hecke algebra and $V_\lambda$ its Specht
module in the Hoefsmit seminormal basis. Write
$R'_j = (T_{j-1}{+}1)\cdots(T_1{+}1)$ and
$B_{\ell+1}^{(\lambda)} V_\lambda := R'_{\ell+1}\cdots R'_n\, V_\lambda$ at
the sharp index $\ell+1$, where $\ell = \ell(\lambda)$. We propose and
test a recursive formula expressing $\dim B_{\ell+1}^{(\lambda)} V_\lambda$
as a sum over length-preserving corners $c_i$ of $\lambda$:
\[
   \dim B_{\ell+1}^{(\lambda)} V_\lambda
   \;=\;
   \sum_{i=1}^{r}\,
   \dim B_{\ell-1}^{(\mu_i)} V_{\mu_i},
   \qquad
   \mu_i := \lambda \setminus \{c_i, c_*(\lambda)\},
\]
where $c_*(\lambda) = (\ell, 1)$ is the column-$1$ bottom and $r$ is the
number of length-preserving corners. The base case ($r = 0$, i.e.\ a hook
or near-hook) is dim~$=1$. Combinatorially, $\dim B_{\ell+1}^{(\lambda)}
V_\lambda$ counts saturated descending chains
$\lambda \supset \lambda^{(1)} \supset \cdots \supset \text{(hook)}$ in
which each step removes the column-$1$ bottom plus one length-preserving
corner.

We verify the formula computationally on nine families covering
$r \in \{1, 2\}$ at $n \in \{8, \ldots, 11\}$, all matching. The
two-corner family is therefore inductively closed: every dimension is
reduced to smaller cases of the same recursion, bottoming out at hooks.

The proof sketch generalises the May-12 three-branch reduction to a
$\binom{r+1}{2}$-branch decomposition. Pairs of length-preserving
corners contribute zero by a $j$-formula leakage argument; pairs of the
form $\{c_i, c_*\}$ contribute the next stage of the recursion;
independence follows from position-of-$n$ and position-of-$(n-1)$
projections on the Hoefsmit seminormal basis.
\end{abstract}

\section{Setup}\label{sec:setup}

\subsection{Notation}

Let $q$ be generic. The \emph{Iwahori--Hecke algebra} $H_q(S_n)$ has
generators $T_1, \ldots, T_{n-1}$ subject to $(T_i - q)(T_i + 1) = 0$ and
the braid relations $T_i T_{i+1} T_i = T_{i+1} T_i T_{i+1}$,
$T_i T_j = T_j T_i$ for $|i-j| \ge 2$. For each $\lambda \vdash n$, write
$V_\lambda$ for the Specht module in the Hoefsmit seminormal basis
$\{v_T : T \in \mathrm{SYT}(\lambda)\}$. The $q$- and $(-1)$-eigenspaces of
$T_j$ are denoted $E_j^+$ and $E_j^-$.

Write
\[
   R'_j \;:=\; (T_{j-1}{+}1)(T_{j-2}{+}1)\cdots(T_1{+}1) \;\in\; H_q(S_j),
\]
and, for $\lambda \vdash n$ with $\ell := \ell(\lambda)$,
\[
   B_{j}^{(\lambda)} V_\lambda \;:=\; R'_j R'_{j+1}\cdots R'_n\, V_\lambda.
\]
We focus on the \emph{sharp} index $j = \ell + 1$:
\[
   B_{\ell+1}^{(\lambda)} V_\lambda \;=\; R'_{\ell+1} R'_{\ell+2}\cdots R'_n\, V_\lambda.
\]

\subsection{Corners}

A cell $c = (r, \lambda_r)$ is a \emph{removable corner} of $\lambda$ if
$\lambda \setminus \{c\}$ is still a partition. We single out two types.

\begin{definition}[Length-preserving corner]
A removable corner $c = (r, \lambda_r)$ is \emph{length-preserving} if
$r < \ell(\lambda)$ and $\lambda_{r+1} < \lambda_r$. Equivalently, removing
$c$ leaves $\ell(\lambda)$ unchanged.
\end{definition}

\begin{definition}[Column-$1$ bottom]
The \emph{column-$1$ bottom} of $\lambda$ is
$c_*(\lambda) := (\ell(\lambda), 1)$. It is always removable (when
$\lambda \ne (1^n)$); removing it decreases $\ell$ by~$1$.
\end{definition}

If $\lambda$ has $r$ length-preserving corners $c_1, \ldots, c_r$, then it
has exactly $r + 1$ removable corners in total: the $c_i$ together with
$c_*$. (For hooks, $r = 0$; for staircases-like shapes, $r$ can be large.)

\subsection{Background}

The starting point is the \emph{rank-zero theorem} (Clio,
2026-05-08--10):
\[
   \Pi^{S_n}|_{V_\lambda} \;=\; 0 \quad \text{whenever } \ell(\lambda) > \lceil n/2 \rceil,
\]
where $\Pi^{S_n} = B_2^{(\lambda)}$. Subsequent work (the
\emph{$j$-formula refinement}, Clio 2026-05-11) showed that at the sharp
index $\ell + 1$,
\[
   B_{\ell+1}^{(\lambda)} V_\lambda \;\subseteq\; E_1^-\!\mid_{V_\lambda},
\]
and is genuinely the bottom non-vanishing stage of the iterated
product. The May-11 \emph{rank-corner formula} conjectured
\[
   \dim B_{\ell+1}^{(\lambda)} V_\lambda \;\stackrel{?}{=}\; \#\{\text{length-pres corners of }\lambda\} \;=\; r.
\]
This was refuted by Clio (2026-05-12, evening~3, paper
\texttt{2026-05-12-evening-4-2-1n-dim3.tex}) at
$\lambda = (4, 2, 1^{n-6})$: the prediction is $r = 2$ but the actual
dimension is $3$. The proof there established a three-branch reduction
yielding $\dim = \dim B^{(\mu_A)} + \dim B^{(\mu_B)} = 2 + 1 = 3$, where
$\mu_A, \mu_B$ are obtained by removing $c_*$ together with the two
length-preserving corners. This suggested a \emph{recursive} replacement
for the rank-corner formula, which we develop here.

\section{The refined formula}\label{sec:formula}

\begin{conjecture}[Refined dimension formula]\label{conj:main}
Let $\lambda \vdash n$ be rank-zero with $\ell = \ell(\lambda)$ and
length-preserving corners $c_1, \ldots, c_r$ ($r \ge 0$), and column-$1$
bottom $c_*(\lambda)$. Then
\begin{equation}\label{eq:main}
   \dim B_{\ell+1}^{(\lambda)} V_\lambda
   \;=\;
   \sum_{i=1}^{r}\,
   \dim B_{\ell-1}^{(\mu_i)} V_{\mu_i},
   \qquad
   \mu_i \;:=\; \lambda \setminus \{c_i, c_*(\lambda)\}.
\end{equation}
The base case ($r = 0$, $\lambda$ a hook) is $\dim = 1$.
\end{conjecture}

Note that $\mu_i$ is a partition of $n - 2$ with
$\ell(\mu_i) = \ell - 1$, so the index $\ell - 1$ on the right is the
\emph{sharp} index $\ell(\mu_i) + 1 - 1 = \ell(\mu_i)$ shifted in the
expected way: the iteration $B_{\ell - 1}^{(\mu_i)}$ on $V_{\mu_i}$
(which lives at $\ell(\mu_i) + 1 = \ell$) is the sharp stage for
$\mu_i$. The recursion therefore closes within the family of rank-zero
shapes.

\begin{remark}[Combinatorial interpretation]\label{rem:chains}
Iterating \eqref{eq:main} expresses $\dim B_{\ell+1}^{(\lambda)} V_\lambda$
as the number of \emph{saturated descending chains}
\[
   \lambda \;=\; \lambda^{(0)} \;\supset\; \lambda^{(1)} \;\supset\; \cdots \;\supset\; \lambda^{(k)} \;=\; (\text{hook}),
\]
where each step removes the column-$1$ bottom of $\lambda^{(t)}$ together
with one of its length-preserving corners:
\[
   \lambda^{(t+1)} \;=\; \lambda^{(t)} \setminus \{c_*(\lambda^{(t)}),\ c_{j_t}^{(t)}\}.
\]
The recursion terminates because each step removes two cells, and a
hook has $r = 0$.
\end{remark}

\begin{example}[Hooks and two-column shapes]\label{ex:hook}
If $\lambda = (k, 1^{n-k})$ with $k \ge 2$, then $\lambda$ has $r = 1$
length-preserving corner (the $(1, k)$ cell when $k \ge 2$), and
$\mu_1 = \lambda \setminus \{(1, k), (n-k+1, 1)\} = (k-1, 1^{n-k-1})$ is
again a hook. The recursion reduces by one column of length one and one
arm-end each step, terminating at $(2, 1)$ or $(1, 1)$ with $\dim = 1$;
hence $\dim B_{\ell+1}^{(\lambda)} V_\lambda = 1$ for every hook.

Similarly, for $\lambda = (2^a, 1^{n-2a})$ with $a \ge 1$, there is a
single length-preserving corner (the $(a, 2)$ cell) and the recursion
descends through the family, again giving $\dim = 1$.
\end{example}

\begin{example}[The $(3, 2, 1^{n-5})$ family: $r = 2$]\label{ex:321}
For $\lambda = (3, 2, 1^{n-5})$, the two length-preserving corners are
$c_1 = (1, 3)$ and $c_2 = (2, 2)$, and $c_* = (n-4, 1)$. The two
partners are
\[
   \mu_1 \;=\; (2, 2, 1^{n-6}),\qquad
   \mu_2 \;=\; (3, 1^{n-5}).
\]
Both are hooks (or hook-like), each contributing $\dim = 1$. Hence the
refined formula predicts $\dim B_{\ell+1}^{(\lambda)} V_\lambda = 1 + 1 = 2$,
agreeing with the May-12 theorem.
\end{example}

\begin{example}[The $(4, 2, 1^{n-6})$ family: $r = 2$, asymmetric]\label{ex:421}
For $\lambda = (4, 2, 1^{n-6})$, corners are $c_1 = (1, 4)$,
$c_2 = (2, 2)$, $c_* = (n-4, 1)$. The two partners are
\[
   \mu_1 \;=\; (3, 2, 1^{n-7}),\qquad
   \mu_2 \;=\; (4, 1^{n-6}).
\]
$\mu_1$ is the $(3, 2, 1^*)$ family from Example~\ref{ex:321}, with
$\dim = 2$; $\mu_2$ is a hook, with $\dim = 1$. The refined formula
predicts $\dim = 2 + 1 = 3$, matching the May-12 evening-3 theorem.
\end{example}

\section{Computational evidence}\label{sec:evidence}

We tested \eqref{eq:main} for nine partition families at
$n \in \{8, 9, 10, 11\}$, working mod $P = 100\,003$ with $q = 11$. The
scripts are at
\texttt{\textasciitilde/projects/scratch/2026-05-13-multi-corner/}.

\begin{center}
\begin{tabular}{lcccc}
\toprule
$\lambda$ & $n$ tested & $r$ & $\dim B_{\ell+1}^{(\lambda)} V_\lambda$ & refined prediction\\
\midrule
Hook $(k, 1^{n-k})$, $k \ge 2$ & various $n$ & $1$ & $1$ & $1$\\
$(2^a, 1^{n-2a})$, $a \ge 1$ & $n \ge 3a$ & $1$ & $1$ & $1$\\
$(3, 2, 1^{n-5})$ & $8, 9$ & $2$ & $2$ & $2 = 1 + 1$\\
$(4, 2, 1^{n-6})$ & $10, 11$ & $2$ & $3$ & $3 = 2 + 1$\\
$(3, 3, 1^{n-6})$ & $10, 11$ & $1$ & $2$ & $2$ (single term, $\mu = (3, 2, 1^*)$, dim $2$)\\
$(4, 3, 1^{n-7})$ & $10, 11$ & $2$ & $5$ & $5 = 2 + 3$\\
$(5, 2, 1^{n-7})$ & $10, 11$ & $2$ & $4$ & $4 = 3 + 1$\\
$(3, 2, 2, 1^{n-7})$ & $10, 11$ & $2$ & $3$ & $3 = 1 + 2$\\
$(4, 2, 2, 1^{n-8})$ & $10, 11$ & $2$ & $6$ & $6 = 3 + 3$\\
$(3, 3, 2, 1^{n-8})$ & $11$ & $2$ & $5$ & $5 = 3 + 2$\\
\bottomrule
\end{tabular}
\end{center}

Every tested case agrees with the refined formula.

\paragraph{An $r = 3$ test.}
We attempted $\lambda = (4, 3, 2, 1^{n-9})$ at $n = 12$ ($r = 3$;
predicted dim $5 + 6 + 5 = 16$). The computation exceeded the $5$-minute
budget on the available machine. Moreover at this $n$ the index
$\ell + 1 = n - 8 = 4$ is at the rank-zero boundary, so the
$j$-formula refinement may not yet apply cleanly. The stable regime
likely begins at $n = 13$, where the predicted value remains $16$. This
test remains open.

\section{Proof sketch (general $r$)}\label{sec:proof}

We outline how the May-12 three-branch reduction generalises. The
argument has five ingredients, all of which are direct generalisations
of pieces established in earlier papers (Clio, May-15, May-19, May-20,
May-12 evening).

\subsection{Ingredient (1): the $(r+1)$-branch reduction}

The Specht module $V_\lambda$ branches under $H_q(S_n) \downarrow H_q(S_{n-1})$
into a direct sum indexed by removable corners:
\[
   V_\lambda \!\downarrow\!_{H_q(S_{n-1})}
   \;\cong\;
   \bigoplus_{\text{removable corner $c$ of }\lambda}
   V_{\lambda \setminus \{c\}}.
\]
With $r + 1$ removable corners (the $c_i$ and $c_*$), there are
$r + 1$ summands. Combining with the $T_{n-1}$-eigenspace decomposition
of $V_\lambda$ into $2$-blocks (indexed by ordered pairs of distinct
corners), one obtains an $H_q(S_{n-2})$-equivariant decomposition
\begin{equation}\label{eq:Eplus-decomp}
   E_{n-1}^+\!\mid_{V_\lambda}
   \;\cong\;
   \bigoplus_{\{c_a, c_b\}}
   V_{\mu_{a,b}},
   \qquad
   \mu_{a,b} \;:=\; \lambda \setminus \{c_a, c_b\},
\end{equation}
where the sum is over unordered pairs of distinct removable corners.
This gives $\binom{r+1}{2}$ branching maps
\[
   \Phi_{a,b}\colon V_{\mu_{a,b}} \;\hookrightarrow\; V_\lambda,
\]
each landing in $E_{n-1}^+$, with mutually disjoint seminormal supports
(distinguished by the positions of $n-1$ and $n$).

This is Theorem~A of the May-20 paper for $r = 2$ (the $(3, 2, 1^*)$
case, $\binom{3}{2} = 3$ branches) and was extended to $r = 3$
($(4, 2, 1^*)$ generates the same structure) in the May-12 evening
paper. The general $r$ statement follows from the standard
Murphy--Hoefsmit theory of $T_{n-1}$-$2$-blocks; we cite it without
proof.

\subsection{Ingredient (2): commutation extends}

The May-15 commutation lemma says that, after applying $R'_n$ to land
in $E_{n-1}^+\!\mid_{V_\lambda}$, the remaining factors
$R'_{n-1}, R'_{n-2}, \ldots, R'_{\ell+1}$ can be pushed past the
branching $\Phi_{a,b}$ at the cost of an outer factor:
\[
   B_{\ell+1}^{(\lambda)} V_\lambda
   \;=\;
   \Big[\prod_{j=\ell - 1}^{n - 2}(T_j {+} 1)\Big]
   \;\cdot\;
   \sum_{\{c_a, c_b\}}
   \Phi_{a, b}\!\left(B_{\ell - 1}^{(\mu_{a,b})} V_{\mu_{a,b}}\right).
\]
The exponent on the outer product reflects that each $R'_k$ for
$\ell + 1 \le k \le n$ contributes a single ``frontier'' factor $(T_{k-1}{+}1)$
after the bulk is absorbed into $\Phi_{a,b}$ by equivariance; the remaining
factors commute through and combine with the corresponding $R'^{(\mu)}_{k-1}$ on
the inside, giving the inner $B_{\ell - 1}^{(\mu_{a, b})}$.

The argument is the same four-step iteration as in the May-12 evening
paper (Theorem~3.1 there), extended to $r + 1$ corners instead of $3$.

\subsection{Ingredient (3): ``bad'' pairs vanish}

\begin{lemma}[Bad-pair vanishing]\label{lem:bad-vanish}
If $\{c_a, c_b\}$ consists of two length-preserving corners (i.e.\
neither is $c_*$), then
\[
   B_{\ell - 1}^{(\mu_{a, b})} V_{\mu_{a, b}} \;=\; 0.
\]
\end{lemma}

\begin{proof}[Proof sketch]
When neither $c_a$ nor $c_b$ is the column-$1$ bottom $c_*$, removing
both does not touch $c_*$: the resulting shape $\mu_{a,b}$ still has its
column-$1$ bottom at row $\ell$ of the original $\lambda$, which equals
$\ell(\mu_{a, b})$ (no rows were deleted because both removed cells were
length-preserving). The sharp index for $\mu_{a, b}$ is therefore
$\ell(\mu_{a, b}) + 1 = \ell + 1$, not $\ell - 1$. The operator
$B_{\ell - 1}^{(\mu_{a,b})}$ applies $R'_{\ell - 1}, \ldots, R'_{\ell + 1}$,
\emph{two extra factors below the sharp index}. By the $j$-formula
refinement (Clio, May-11), the iteration past the sharp index forces
descent through $E_1^-$, where the rightmost $(T_1{+}1)$ in each $R'_k$
annihilates the eigenspace. (This is the same $E_1^-$-leakage mechanism
that proved Theorem~3.3 of the May-12 evening paper.)
\end{proof}

\subsection{Ingredient (4): ``good'' pairs contribute the recursion}

\begin{lemma}[Good-pair matching]\label{lem:good-match}
If $\{c_a, c_b\} = \{c_i, c_*\}$ for some length-preserving corner $c_i$,
then
\[
   \mu_{i, *} \;=\; \mu_i \;=\; \lambda \setminus \{c_i, c_*\}
\]
has $\ell(\mu_i) = \ell - 1$, and the sharp index for $\mu_i$ is
$\ell(\mu_i) + 1 = \ell$. The branching contribution at $\{c_i, c_*\}$
is then
\[
   \Phi_{i, *}\!\left(B_{\ell - 1}^{(\mu_i)} V_{\mu_i}\right)
   \;=\;
   \Phi_{i, *}\!\left(B_{\ell(\mu_i) + 1 - 1}^{(\mu_i)} V_{\mu_i}\right),
\]
which is exactly the sharp-stage $B_{\ell(\mu_i)}^{(\mu_i)} V_{\mu_i}$
mapped through $\Phi_{i, *}$ — and \emph{this is the recursion's next stage
on $\mu_i$ at one step before the sharp index}, i.e.\
$B_{\ell(\mu_i) + 1}^{(\mu_i)} V_{\mu_i}$ before pulling back through $R'_{\ell(\mu_i) + 1}$.

In particular, applying the outer factors $\prod_{j = \ell - 1}^{n - 2}(T_j{+}1)$ and
the inner $R'^{(\mu_i)}_{\ell(\mu_i) + 1}$ together reconstructs the full
sharp-stage object on $\mu_i$.
\end{lemma}

\noindent The point is that the index shift $\ell - 1 = \ell(\mu_i)$ on the
right of \eqref{eq:main} is \emph{exactly} the sharp index for $\mu_i$.
The recursion therefore closes: each good-pair contribution is itself a
sharp-stage object for the smaller partition $\mu_i$.

\subsection{Ingredient (5): independence (lower bound)}

\begin{proposition}[Independence of good branches]\label{prop:indep}
The $r$ good branches $\Phi_{i, *}(B_{\ell-1}^{(\mu_i)} V_{\mu_i})$,
$1 \le i \le r$, inject independently into $V_\lambda$. Hence
$\dim B_{\ell+1}^{(\lambda)} V_\lambda \;\ge\; \sum_i \dim B_{\ell-1}^{(\mu_i)} V_{\mu_i}$.
\end{proposition}

\begin{proof}[Proof sketch]
The position-of-$n$ projection $\pi_n^{(c)}: V_\lambda \to V_\lambda$
onto SYTs with letter $n$ at cell $c$ commutes with $T_j{+}1$ for
$j \le n - 2$ (Clio, May-12, Lemma~4.1 — letters $\le n - 1$ cannot
move letter $n$). Each $\Phi_{i, *}$ has $\{n - 1, n\}$ supported on
$\{c_i, c_*\}$, so $\pi_n^{(c_j)} \Phi_{i, *} = 0$ for $j \ne i$. Hence
projection onto $\pi_n^{(c_i)}$ separates the $i$th branch from all
other branches at the level of seminormal supports.

Within a single $i$-branch, independence of the different elements of
$B_{\ell - 1}^{(\mu_i)} V_{\mu_i}$ is inherited from the recursive
hypothesis on $\mu_i$ by injectivity of $\Phi_{i, *}$; a second
position-of-$(n-1)$ projection identifies which sub-branch of $\mu_i$
(i.e.\ which good pair inside $\mu_i$) each summand comes from, as in
the May-12 evening paper's Lemma~5.2.
\end{proof}

\subsection{Putting it together}

\begin{theorem}[Recursion, conditional]\label{thm:main}
Subject to ingredients (1)--(5), the refined formula
\eqref{eq:main} holds for every rank-zero $\lambda$ with
$r \ge 1$ length-preserving corners. Iterating, the dimension equals
the number of saturated descending chains of Remark~\ref{rem:chains}.
\end{theorem}

\begin{proof}[Proof]
Combine ingredients (1)--(3) for the upper bound: the
$\binom{r+1}{2}$-branch decomposition contains $\binom{r}{2}$ bad-pair
contributions (length-preserving + length-preserving), which vanish by
Lemma~\ref{lem:bad-vanish}, and $r$ good-pair contributions
$\{c_i, c_*\}$, which map injectively. Each good contribution has the
inner dimension $\dim B_{\ell - 1}^{(\mu_i)} V_{\mu_i}$ by
Lemma~\ref{lem:good-match}. The supports of the good branches are
pairwise disjoint (Ingredient~(1)). Hence
\[
   \dim B_{\ell + 1}^{(\lambda)} V_\lambda
   \;\le\;
   \sum_{i = 1}^{r}
   \dim B_{\ell - 1}^{(\mu_i)} V_{\mu_i}.
\]
The lower bound is Proposition~\ref{prop:indep}. Equality follows.
\end{proof}

\subsection{Modulo what?}

The proof is rigorous modulo two technical inputs:

\begin{enumerate}[leftmargin=2em,topsep=0.3em,itemsep=0.2em,label=(\roman*)]
\item \textbf{Phase A generalisation at general $r$.} The $T_{n-1}$-$2$-block
decomposition \eqref{eq:Eplus-decomp} for $r + 1$ corners
must give exactly $\binom{r+1}{2}$ pieces with the stated branching
maps. This is standard in Murphy--Hoefsmit theory and follows from the
seminormal eigenvalue structure (axial-distance contents at distinct
corner pairs); we cite it without proof.

\item \textbf{Hoefsmit nonvanishing checks.} The position-of-$n$ and
position-of-$(n-1)$ projections used in Proposition~\ref{prop:indep}
must give nonzero output at each step. Each check is a \emph{finite}
case analysis on a $2$-block at known axial distance, which is
non-degenerate for the rank-zero regime $n \ge n_0(\lambda)$. These
have been verified computationally for all nine tested shapes; for
each individual family they reduce to a finite enumeration analogous
to Lemma~4.2 of the May-12 evening paper.
\end{enumerate}

\section{Implications and open questions}\label{sec:open}

\paragraph{The $r = 2$ family is inductively closed.}
For every $\lambda$ with $r = 2$ length-preserving corners, the formula
\eqref{eq:main} reduces $\dim B_{\ell + 1}^{(\lambda)} V_\lambda$ to a
sum of two terms, each of which is the dimension at a shape with
$r' \in \{1, 2\}$ at smaller $n$. The $r' = 1$ shapes are hooks and
two-column shapes (Example~\ref{ex:hook}) with dim $= 1$. The $r' = 2$
shapes recurse, and at most $\lfloor (n - 2)/2 \rfloor$ levels of
recursion suffice (each step removes two cells). So the entire $r = 2$
family is reduced, modulo the technical ingredients above, to a
closed-form computation.

\paragraph{$r = 3$ awaits.}
The smallest $r = 3$ shape is $\lambda = (4, 3, 2, 1^{n - 9})$. The
refined formula predicts dim $= 5 + 6 + 5 = 16$ from the three good
pairs $(c_1, c_*), (c_2, c_*), (c_3, c_*)$, which give
$\mu_1 = (3, 3, 2, 1^*)$ (dim $5$), $\mu_2 = (4, 2, 2, 1^*)$ (dim $6$),
$\mu_3 = (4, 3, 1^*)$ (dim $5$). We could not compute this at $n = 12$
in our budget; we expect $n = 13$ to be both more tractable (well
inside the rank-zero stable regime) and decisive for the prediction.

\paragraph{Closed form?}
The recursion identifies $\dim B_{\ell + 1}^{(\lambda)} V_\lambda$ with
the number of saturated descending chains in a corner-pair poset on
$\lambda$. It is natural to ask:
\begin{itemize}[leftmargin=2em,topsep=0.3em,itemsep=0.2em]
\item Is there a closed-form expression, e.g.\ a product over
local corner data of $\lambda$?
\item Does the chain count match a known combinatorial statistic
(e.g.\ a count of certain SYTs, paths in Young's lattice, or
$0$-cells in a derived complex)?
\item Is the chain length (i.e.\ how many steps until $\lambda$
becomes a hook) given by $\sum_i \sigma(c_i)$ for some local
statistic $\sigma$? Each step strips $c_*$ plus one length-pres
corner, so the total number of steps is one less than $\ell - 1$
plus the number of length-pres corners stripped along the way —
which suggests $\sigma$ should track ``distance from being a
hook''.
\end{itemize}

\paragraph{Status.} The formula \eqref{eq:main} matches the
computational data on every tested family ($r \in \{1, 2\}$,
$n \le 11$). The general proof is structurally complete modulo
ingredients~(i)--(ii) of Section~\ref{sec:proof}. The next concrete
target is the $r = 3$ verification at $n = 13$.

\end{document}
